
\documentclass{article}
\usepackage{colortbl}
\usepackage{makecell}
\usepackage{multirow}
\usepackage{supertabular}

\begin{document}

\newcounter{utterance}

\centering \large Interaction Transcript for game `cladder', experiment `full\_v1.5\_default', episode 916 with refineonly{-}llama{-}3.1{-}8b{-}instruct.
\vspace{24pt}

{ \footnotesize  \setcounter{utterance}{1}
\setlength{\tabcolsep}{0pt}
\begin{supertabular}{c@{$\;$}|p{.15\linewidth}@{}p{.15\linewidth}p{.15\linewidth}p{.15\linewidth}p{.15\linewidth}p{.15\linewidth}}
    \# & \multicolumn{2}{c}{Player} && \multicolumn{2}{c}{Game Master} \\
    \hline

    \theutterance \stepcounter{utterance}  
    & & & \multicolumn{4}{p{0.6\linewidth}}{
        \cellcolor[rgb]{0.9,0.9,0.9}{
            \makecell[{{p{\linewidth}}}]{
                \texttt{\tiny{[P1$\langle$GM]}}
                \texttt{You are an expert in causal inference. The following question is not a typical commonsense query, but rather a meticulously designed question created by a professor specializing in causal inference, intended to assess the students' mastery of the course content.} \\
\texttt{Imagine a self{-}contained, hypothetical world with only the following conditions, and without any unmentioned factors or causal relationships: Smoking has a direct effect on effort and college admission. Effort has a direct effect on college admission. For nonsmokers and are lazy, the probability of college admission is 47\%. For nonsmokers and are hard{-}working, the probability of college admission is 88\%. For smokers and are lazy, the probability of college admission is 23\%. For smokers and are hard{-}working, the probability of college admission is 54\%. For nonsmokers, the probability of being hard{-}working is 40\%. For smokers, the probability of being hard{-}working is 61\%. If we disregard the mediation effect through effort, would smoking negatively affect college admission?} \\
\texttt{Start your answer with 'Yes' or 'No', followed by additional reasoning or evidence to support your explanation.} \\
            }
        }
    }
     \\ \\

    \theutterance \stepcounter{utterance}  
    & \multicolumn{4}{p{0.6\linewidth}}{
        \cellcolor[rgb]{0.561,0.737,0.561}{
            \makecell[{{p{\linewidth}}}]{
                \texttt{\tiny{[P1$\rangle$GM]}}
                \texttt{No.} \\
\\ 
\texttt{The average treatment effect (ATE) is calculated as ATE = P(Y(1)) {-} P(Y(0)). We can calculate P(Y(1)) and P(Y(0)) using the law of total probability. } \\
\\ 
\texttt{For smokers:} \\
\texttt{P(Y(1)) = P(Y(1)|H)P(H|S) + P(Y(1)|L)P(L|S)} \\
\texttt{= 54\% * 61\% + 23\% * 39\%} \\
\texttt{= 0.54 * 0.61 + 0.23 * 0.39} \\
\texttt{= 0.3303 + 0.0897} \\
\texttt{= 0.420} \\
\\ 
\texttt{For nonsmokers:} \\
\texttt{P(Y(0)) = P(Y(0)|H)P(H|S=0) + P(Y(0)|L)P(L|S=0)} \\
\texttt{= 88\% * 40\% + 47\% * 60\%} \\
\texttt{= 0.88 * 0.40 + 0.47 * 0.60} \\
\texttt{= 0.352 + 0.282} \\
\texttt{= 0.634} \\
\\ 
\texttt{Now, we can calculate the ATE:} \\
\texttt{ATE = P(Y(1)) {-} P(Y(0))} \\
\texttt{= 0.420 {-} 0.634} \\
\texttt{= {-}0.214} \\
\\ 
\texttt{Since the ATE is negative, smoking would negatively affect college admission if we disregard the mediation effect through effort.} \\
            }
        }
    }
    & & \\ \\

    \theutterance \stepcounter{utterance}  
    & & & \multicolumn{2}{p{0.3\linewidth}}{
        \cellcolor[rgb]{0.9,0.9,0.9}{
            \makecell[{{p{\linewidth}}}]{
                \texttt{\tiny{[GM$|$GM]}}
                \texttt{No} \\
            }
        }
    }
    & & \\ \\

    \theutterance \stepcounter{utterance}  
    & & & \multicolumn{2}{p{0.3\linewidth}}{
        \cellcolor[rgb]{0.9,0.9,0.9}{
            \makecell[{{p{\linewidth}}}]{
                \texttt{\tiny{[GM$|$GM]}}
                \texttt{yes} \\
            }
        }
    }
    & & \\ \\

    \theutterance \stepcounter{utterance}  
    & & & \multicolumn{2}{p{0.3\linewidth}}{
        \cellcolor[rgb]{0.9,0.9,0.9}{
            \makecell[{{p{\linewidth}}}]{
                \texttt{\tiny{[GM$|$GM]}}
                \texttt{game\_result = LOSE} \\
            }
        }
    }
    & & \\ \\

\end{supertabular}
}

\end{document}
